We deduce a topological and homotopical Lefschetz fixed point 
via sheaves and local systems. 

\subsection{Topological Lefschetz}

\begin{defn}
    Let \mdef{$\Top$} denote the category of \tdef{locally compact Hausdorff spaces}. For any $\base \in \Top$ we write 
    $\mdef{\pTop{\base}}\subseteq \Top_{/\base}$ 
    for the full subcategory spanned by the \tdef{proper} maps $p\colon X\to \base$ respectively. Moreover, for any $X \in \Top$ we denote 
\[\mdef{\pSp{X}}\coloneq \Shv(X;\Sp) \in \CAlg(\Prlst)\]
for the stably presentably monoidal $\infty$-category of \tdef{sheaves of spectra} on $X$ with unit $\SS_{X}\in \pSp{X}$.
\end{defn}

Throughout the rest of this section, we fix a base space $M\in \Top$.

\begin{cons}
Consider the geometric context $(\pTop{\base},\all)$, let
$\cI$ denote the class of open immersions and $\cP$ the class of proper maps.

By \todo{cite Volpe?}, the assignment 
$X\mapsto \pSp{X}$ refines to a symmetric monoidal functor $(\infty,2)$-functor $\sSpan{\Top}_{\cP,\cI}\to \Prlstf$. Pre-composing
with the functor $\sSpan{\pTop{\base}}_{\cP,\cI} \to \sSpan{\Top}_{\cP,\cI}$ that forgets the structure map, we obtain an (a priori) 
lax symmetric monoidal functor $\pSpf{(-)}\colon \sSpan{\pTop{\base}}_{\cP,\cI}\too \Prlstf$.

Now note that by the proper base change formula \todo{cite} this functor
is in fact strong monoidal. Moreover, by general nonsense, we obtain
a lift to a symmetric monoidal $(\infty,2)$-functor 
\[ \mdef{\pSpf{(-)}}\colon \sSpan{\pTop{\base}}_{\cP,\cI}\too \pPrlf{\base}\,,\]
which we shall use throughout this section.
\end{cons}

Let us unwind the data of this $(\infty,2)$-functor that will
be relevant for us.

\begin{notation}
For any morphism $f\colon X\to Y \in \pTop{\base}$ we write 
\[ \mdef{f^*}\colon \pSp{Y}\to \pSp{X}\quad \text{and}\quad  \mdef{f_!} \colon \pSp{X}\to \pSp{Y}\]
for the image under the functor $\pSpf{(-)}$ of the spans 
\[ \left( Y \xleftarrow{f} X \xrightarrow{\id}X\right) \quad \text{and} \quad  \left(X \xleftarrow{\id} X \xrightarrow{f} Y \right)\,, \]
respectively.
For any $(X,p)\in \pTop{\base}$ we write
\[ \mdef{\cptSec{X}[\base]}\coloneq p_!\SS_{X} \in \pSp{\base}\,. \]
for the \tdef{relative cohomology} of $X$.

Note that, $\cptSec{X}[\base]$ is equivalently the image
of the span $ \base \xleftarrow{p} X \xrightarrow{p} \base$
under the functor 
\[\Omega \pSpf{(-)}\colon \Omega \sSpan{\pTop{\base}}_{\cP,\cI} \to
\Omega \pPrlf{\base}\simeq \pSp{\base}\,.\]
Moreover, for any proper map $f\colon (X,p)\to (Y,q)$ in $\Top{/\base}$ use
the notation
\[ \mdef{\pb{f}}\colon \cptSec{Y}[\base]\to \cptSec{X}[\base] \]
for the pullback on relative compactly supported cohomology, i.e.~the image of the 2-morphism 
\[
\begin{tikzcd}
	& X & \\
	\base & X & \base \\
	& Y
	\arrow["p"', from=1-2, to=2-1]
	\arrow["p", from=1-2, to=2-3]
	\arrow["\id"', from=2-2, to=1-2]
	\arrow["f", from=2-2, to=3-2]
	\arrow["q", from=3-2, to=2-1]
	\arrow["q"', from=3-2, to=2-3]
\end{tikzcd}
\]
under the $(\infty,2)$-functor $\pSpf{(-)}\colon \sSpan{\pTop{\base}}_{\cP,\cI} \to \pPrlf{\base}$.
\end{notation}


\begin{lem}\label{lem: top trace identification}
    Let $f\colon (X,p)\to (X,p)$ be a self map in $\pTop{\base}$ with
    fixed points $i\colon X^f\to X$. Then we
    have an equivalence:
    \[ \tr{f_!\colon \pSp{X}\to \pSp{X}}{\pPrl{\base}}\simeq \cptSec{X^f}[\base]\in \pSp{\base}\,.\]
\end{lem}
\begin{proof}
  By the naturality statment of \cref{thm: trace functor}, the symmetric monoidal 
  $(\infty,2)$-functor $\pSpf{(-)}$ induces a symmetric monoidal functor 
  \[ \pSpf{(-)}\colon \Endtrl(\sSpan{\pTop{\base}}_{\cP,\cI}) 
  \too \Endtrl(\pPrlf{\base})\,,\]
  which is compatible with trace. Thus, we learn that 
 the trace of the functor $f_!\colon \pSp{X}\to \pSp{X}$
 is computed as the image of the trace of the span 
 $(X\xleftarrow{\id} X \xrightarrow{f}X)$, which by \cref{lem: trace in Span}
 is given by the fixed points $\base \leftarrow X^f \to \base$,
 which was the claim. 
\end{proof}

\begin{cons}\label{cons: Lefschetz index}
    Let $f\colon (X,p)\to (X,p)$ be a self map in $\pTop{\base}$ and
    assume that the projection $p\colon X\to M$ is a shape submersion. Consider the following diagram in $\pPrlf{\base}$:
   \[\begin{tikzcd}
	{\pSp{X}} & {\pSp{\base}} \\
	{\pSp{X}} & {\pSp{\base}}
	\arrow["{p_!}", from=1-1, to=1-2]
	\arrow["{f_!}"', from=1-1, to=2-1]
	\arrow["\id", from=1-2, to=2-2]
	\arrow[Rightarrow, from=2-1, to=1-2]
	\arrow["{p_!}"', from=2-1, to=2-2]
\end{tikzcd} \]
where the 2-cell is induced by the equation
$pf=p$. Since $p$ is a proper shape submersion,
the functor $p_!$ admits an internal right adjoint by REF/CITE,
and hence this defines a morphism $(\pSp{X},f_!)\to (\pSp{\base},\id)$ in $\Endtrl(\pPrlf{\base})$. Applying the symmetric
monoidal functor 
\[ \TrcL\colon \Endtrl(\pPrlf{\base}) \to \pSp{\base}\]
of \cref{thm: trace functor}
and using the identification of \cref{lem: top trace identification} we obtain
a map:
\[ \mdef{\lambda^f}\colon \cptSec{X^f}[\base]
\to \SS_{\base} \in \pSp{\base}\]
which we call the \tdef{Lefschetz index} of $f$. 
    \end{cons}

The following technical lemma will allow us to identify the pullback
of the Lefschetz index to the point.

\begin{lem}\label{lem: proper internal right adj}
    For any $(X,p)\in \pTop{\base}$ the span 
    $(\base \xleftarrow{p} X \xrightarrow{\id} X) \in \sSpan{\pTop{\base}}_{\cP,\cI}$ admits an internal
    right adjoint given by $(X \xleftarrow{\id} X \xrightarrow{p} \base)$.
\end{lem}
\begin{proof}
    This is a special case of \cref{lem: internal adjoints}. 
\end{proof}

We can now prove the main result of this section:

\begin{prop}\label{prop: Top Lefschetz}
Let $f\colon (X,p)\to (X,p)$ be a self map in $\pTop{\base}$ and assume
that $p\colon X\to \base$ is a shape submersion. 
Let $\lambda^f\colon \cptSec{X^f}[\base] \to \SS_{\base}$ be
the Lefschetz index of \cref{cons: Lefschetz index} and write 
$q\colon X^f\to \base$ for the terminal map.
The composite 
\[ \SS_{\base} \xrightarrow{\pb{q}} \cptSec{X^f}[\base]
\xrightarrow{\lambda^f} \SS_{\base}\,\]
is given by the trace of the pullback map 
\[ \pb{f}\colon \cptSec{X}[\base]\to \cptSec{X}[\base] \]
on relative cohomology.
\end{prop}
\begin{proof}
    Consider the commutative diagram in $\Endtrl(\pPrlf{\base})$:
\begin{equation}\label{eq: categorical factorization top}
   \begin{tikzcd}
	{\pSp{\base}} & {\pSp{X}} & {\pSp{\base}} \\
	{\pSp{\base}} & {\pSp{X}} & {\pSp{\base}}
	\arrow["{p^*}", from=1-1, to=1-2]
	\arrow["\id"', from=1-1, to=2-1]
	\arrow["{p_!}", from=1-2, to=1-3]
	\arrow["{f_!}", from=1-2, to=2-2]
	\arrow["\id", from=1-3, to=2-3]
	\arrow["\alpha", Rightarrow, from=2-1, to=1-2]
	\arrow["{p^*}"', from=2-1, to=2-2]
	\arrow["\beta", Rightarrow, from=2-2, to=1-3]
	\arrow["{p_!}"', from=2-2, to=2-3]
\end{tikzcd} 
    \end{equation}
    where $\alpha\colon f_!p^*\to p^*$ is adjoint to the preferred equivalence $f^*p^* = p^*$ and $\beta$ is the preferred
    equivalence $p_!f_!=p_!$. Applying the functor 
    \[ \TrcL\colon \Endtrl(\pPrlf{\base})\to \pSp{\base}\,\]
    we obtain a commutative diagram in $\pSp{\base}$ of
    the form:
    \[
\begin{tikzcd}
	{} & {\cptSec{X^f}[\base]} & \\
	{\SS_{\base}} && {\SS_{\base}}
	\arrow["{\lambda^f}", from=1-2, to=2-3]
	\arrow["{\TrcL(p^*,\alpha)}", from=2-1, to=1-2]
	\arrow["{\TrcL(p_!p^*, \beta\alpha)}"', from=2-1, to=2-3]
\end{tikzcd}
    \]
    The map $\beta\alpha\colon p_!p^*\SS_{\base}\to p_!p^*\SS_{\base}$ is by definition the pullback $\pb{f}$
    along $f$ and so by construction of the functor 
    $\TrcL$ REF, we have that
    \[ \TrcL(p_!p^*,\beta\alpha)= \tr{\pb{f}}{\pSp{\base}}\,.\]
    It remains to show that $\TrcL(p^*,\alpha)$ is given
    by the pullback $\pb{q}$ along the projection 
    map $q\colon X^{f}\to \base$.
    Indeed, by \cref{lem: proper internal right adj} the span 
    $(\base \xleftarrow{p} X \xrightarrow{\id} X)$ admits an internal
    right adjoint and thus defines a map:
    \[
    \begin{tikzcd}
	\base & X & X \\
	\base && X \\
	\base & X & X 
	\arrow["p"', from=1-2, to=1-1]
	\arrow["\id", from=1-2, to=1-3]
	\arrow["\id"', from=2-1, to=1-1]
	\arrow["\id", from=2-1, to=3-1]
	\arrow["f", from=2-3, to=1-3]
	\arrow["\id"', from=2-3, to=3-3]
	\arrow["p"', from=3-2, to=3-1]
	\arrow["\id", from=3-2, to=3-3]
\end{tikzcd}
\]
in $\Endtrl(\sSpan{\pTop{\base}}_{\cP,\cI})$ whose image under the functor 
$\pSp{(-)}$ is the left hand square of \eqref{eq: categorical factorization top}. Thus, unwinding the definitions and using the naturality of the
functor $\Trc$, we learn that the map $\TrcL(p^*,\alpha)$
is the image of the 2-morphism: 
\[
\begin{tikzcd}
	& {X^f} & \\
	\base & {X^f} & \base \\
	& \base
	\arrow["q"', from=1-2, to=2-1]
	\arrow["q", from=1-2, to=2-3]
	\arrow["\id"', from=2-2, to=1-2]
	\arrow["q", from=2-2, to=3-2]
	\arrow[from=3-2, to=2-1]
	\arrow[from=3-2, to=2-3]
\end{tikzcd}
\]
under the functor $\Omega \pSpf{(-)}$, which is exactly the pullback along 
$q\colon X^f \to \base$. 
\end{proof}

\subsection{Homotopical Lefschetz}\label{subsec: homotopical lefschetz}
\begin{defn}
    For some $\base \in \Top$ we denote the $\infty$-category
    of sheaves of spaces on $\base$ by:
    \[\mdef{\pSpc{\base}}\coloneq \Shv(\base,\Spc) \in \Prl\,.\]
    we denote the terminal object
    of $\pSpc{\base}$ by $\pt_{\base}$. 
    Moreover. for any $A\in \pSpc{\base}$ we define:
    \[ \mdef{\pSp{\base}[A]}\coloneq \Shv((\pSpc{\base})_{/A};\Sp)\in \CAlg(\Prlst)\,, \]
    and write $\SS_{A}$ for the monoidal unit of $\pSp{\base}[A]$.
\end{defn}

Throughout this section, we fix a hypercomplete base space $M\in \Top$.
By extending the functor $\pSp{\base}[(-)]$ to symmetric monoidal 
$(\infty,2)$-categorical six functor formalism.

\begin{cons}
Consider $\pSp{\base}[(-)]$ as a functor $\pSpc{\base}^{\op}
\to \pPrl{\base}$ via pullback of sheaves and the pointwise
$\pSp{\base}$-linear structure induced by the terminal maps 
and the natural identification $\pSp{\base}[\pt_{\base}]\simeq \pSp{\base}$. Write $\enquote{\all}$ for the class of 
all maps and $\enquote{\simeq}$ for the class
of invertible maps in $\pSpc{\base}$. It is clear that 
$(\pSpc{\base}, \all)$ is a geometric context and moreover that the
pair $(\all, \simeq)$ forms a suitable decomposition in the sense
of \cite{Mann2022}. 

We further claim that $\pSp{\base}[(-)]$ is 
$(\all, \simeq)$-biadjointable in the sense of~\cite[Definition 2.4]{Cnossen_Lenz_Linskens2026}. 
This amounts to checking that series of naturally defined Beck-Chevalley
transformations are equivalences. Since we assume that $\base$ is 
hypercomplete, the topos $\pSpc{\base}$ is hypercomplete as well,
so we may verify this on stalks. Using the \'etale base change 
formula \todo{cite} for any map $x\colon \pt \to \base$, we reduce
to the case $\base = \pt$, where the claim follows from \todo{Find citation}. \todo{Expand this out a little bit}

Thus, by \cite[Theorem A]{Cnossen_Lenz_Linskens2026} the assignment
$A\mapsto \pSp{\base}[A]$ admits a natural refinement to a symmetric
monoidal \todo{Give an argument for why it's not just lax}
$(\infty,2)$-functor:
\[ \mdef{\pSpf{\base}[(-)]}\colon \sSpan{\pSpc{\base}}_{\simeq,\all} \too \pPrlf{\base}\,.\]
\end{cons}

As before, let us unwind the data of this $(\infty,2)$-functor that
is relevant for us.

\begin{notation}
For any morphism $f\colon A\to B \in \pSpc{\base}$ we write 
\[ \mdef{f^*}\colon \pSp{\base}[B]\to \pSp{\base}[A]\quad \text{and}\quad  \mdef{f_!} \colon \pSp{\base}[A]\to \pSp{\base}[B]\]
for the image of the spans 
\[ \left( B \xleftarrow{f} A \xrightarrow{\id}A \right) \quad \text{and} \quad  \left(A \xleftarrow{\id} A \xrightarrow{f} B \right)\,, \]
respectively.
For any $A\in \pSpc{\base}$ with terminal map $p\colon A\to \pt_{\base}$ we denote
by
\[\mdef{\SS_{\base}[A]}\coloneq p_!\SS_{X} \in \pSp{\base}\]
the \tdef{relative spherical homology} of $A$. Note that, this
is equivalently the image
of the span $\pt_\base\leftarrow A \rightarrow \pt_\base$ under
the functor
\[\Omega\pSpf{\base}[(-)]\colon \sSpan{\pSpc{\base}}_{\simeq,\all}\too \pPrlf{\base}\,.\] 
Moreover, for any map 
$f\colon A\to B\in \pSpc{\base}$ we write 
\[\mdef{\tf{f}}\colon \SS_{\base}[A]\to \SS_{\base}[B]\]
for the \tdef{transfer map} induced from the 2-morphism:
\[\begin{tikzcd}
	& B & \\
	\pt_\base & A & \pt_\base \\
	& A
	\arrow[from=1-2, to=2-1]
	\arrow[from=1-2, to=2-3]
	\arrow["f"', from=2-2, to=1-2]
	\arrow["\id", from=2-2, to=3-2]
	\arrow[from=3-2, to=2-1]
	\arrow[from=3-2, to=2-3]
\end{tikzcd}\]
via the functor $\Omega\pSpf{\base}[(-)]$.
\end{notation}


\begin{lem}\label{lem: etale internal right adj}
For any $A\in \pSpc{\base}$the span 
$(A\xleftarrow{\id} A \to \pt_\base) \in \sSpan{\pSpc{\base}}_{\simeq, \all}$ admits an internal right adjoint given by 
$(\pt_\base \leftarrow A \xrightarrow{\id} A)$.
\end{lem}
\begin{proof}
This is a special case of \cref{lem: internal adjoints}. 
\end{proof}

\begin{lem}\label{lem: homo trace identification}
    Let $f\colon A\to A$ be a self map in $\pSpc{\base}$ with homotopy
    fixed points $i\colon \cL^fA \to A$. We have an equivalence 
    \[ \tr{f^* \colon \pSp{\base}[A]\to \pSp{\base}[A]}{\pPrl{\base}}
    \simeq \SS_{\base}[\cL^fA] \in \pSp{\base}\,
    \]
   which is natural in spans $A\xleftarrow{i} C \xrightarrow{g} B$ where
   $i$ is an equivalence and $g$ is any map.
\end{lem}
\begin{proof}
    Completely analogous to \cref{lem: top trace identification} making
    use of the general calculation of traces in span categories done in \cref{lem: trace in Span}.
\end{proof}


\begin{cons}\label{cons: Reidemeister trace}
   Let $f\colon A\to A$ be a self map in $\pSpc{\base}$ and
   write $p\colon A\to \pt_\base$ for the terminal map.
   Consider the following diagram in $\pPrlf{\base}$:
   \[\begin{tikzcd}
	{\pSp{\base}} & {\pSp{\base}[A]} \\
	{\pSp{\base}} & {\pSp{\base}[A]}
	\arrow["{p^{*}}", from=1-1, to=1-2]
	\arrow["{\id}"', from=1-1, to=2-1]
	\arrow["f^*", from=1-2, to=2-2]
	\arrow[Rightarrow, from=2-1, to=1-2]
	\arrow["{p^*}"', from=2-1, to=2-2]
\end{tikzcd} \]
where the 2-cell is induced by the natural identification
$pf=p$. Since any map in $\Spc_{\base}$ is \'etale \todo{REF/CITE} this 
datum defines a map
\[ (\pSp{\base},\id)\to (\pSp{\base}[A],f^*)\]
in $\Endtrl(\pPrlf{\base})$. Applying
the symmetric monoidal functor $\TrcL$ 
and using the identification of \cref{lem: homo trace identification}, we obtain a map:
\[ \mdef{\rho_f}\colon \SS_{\base} \to \SS_{\base}[\cL^fA] \]
which we call the \tdef{Reidemeister trace} of $f$. 
\end{cons}

\begin{prop}
    Let $f\colon A\to A$ be a self map in $\pSpc{\base}$ \todo{add assumption}with
    terminal map $p\colon A\to \pt_\base$,
   homotopy fixed points $i\colon \cL^fA \to A$
    and Reidemeister
    trace $\rho_f\colon \SS_{\base}\to \SS_{\base}[\cL^f A]$.
    The composite
    \[ \SS_{\base}\xrightarrow{\rho_f} \SS_{\base}[\cL^f A] \xrightarrow{\tf{ip}} \SS_{\base}\]
    is given by the trace of the transfer 
    \[ \tf{f}\colon \SS_{\base}[A]\to \SS_{\base}[A]\]
    on relative spherical homology.
\end{prop}
\begin{proof}
    The proof is completely analogous to the one of \cref{prop: Top Lefschetz}. Indeed, consider the commutative diagram in $\Endtrl(\pPrlf{\base})$:
    \[
   \begin{tikzcd}
	{\pSp{\base}} & {\pSp{\base}[A]} & {\pSp{\base}} \\
	{\pSp{\base}} & {\pSp{\base}[A]} & {\pSp{\base}}
	\arrow["{p^*}", from=1-1, to=1-2]
	\arrow["\id"', from=1-1, to=2-1]
	\arrow["{p_\sharp}", from=1-2, to=1-3]
	\arrow["{f^*}", from=1-2, to=2-2]
	\arrow["\id", from=1-3, to=2-3]
	\arrow["\gamma", Rightarrow, from=2-1, to=1-2]
	\arrow["{p^*}"', from=2-1, to=2-2]
	\arrow["\delta", Rightarrow, from=2-2, to=1-3]
	\arrow["{p_\sharp}"', from=2-2, to=2-3]
\end{tikzcd} 
    \]
where $\gamma$ is the preferred equivalence $f^{*}p^*=p^*$ and 
$\delta$ is adjoint to the preferred equivalence $p_\sharp f_{\sharp}= p_{\sharp}$. Applying the functor $\TrcL$, we obtain a commutative diagram:
    \[\begin{tikzcd}
	{} & {\SS_{\base}[\cL^fA]} & \\
	{\SS_{\base}} && {\SS_{\base}}
	\arrow["{\TrcL(p_{\sharp},\delta})", from=1-2, to=2-3]
	\arrow["{\rho_f}", from=2-1, to=1-2]
	\arrow["{\TrcL(p_\sharp p^*, \delta\gamma)}"', from=2-1, to=2-3]
\end{tikzcd}\]
   Unwinding the definitions the lower horizontal map naturally identifies with the trace
   of the pushforward map. Moreover, we can identify the right hand
   diagnoal map with the pushforward along the terminal map 
   $ip\colon \cL^fA\to \pt_\base$ using \cref{lem: etale internal right adj}.
\end{proof}
